\chapter{Заключение}


В ходе выполнения работы получены следующие основные результаты.

Изучены основные концепции кватернионов Гамильтона и бикватернионов, которые применимы к описанию пространственных преобразований. Рассмотрены алгебра кватернионов в нотации однородных координат, операции сопряжения, скалярное и векторное произведения, а также представление единичного кватерниона в тригонометрической форме. На этой основе изложены формулы вращения точки и вектора вокруг оси, проходящей через начало координат, с помощью сэндвич-формулы.

Рассмотрено определение дуальных бикватернионов тремя способами: через таблицу умножения базисных элементов, через процедуру удвоения Кейли–Диксона и через дуальное представление кватернионных формул. Введены операции над бикватернионами, включая три типа сопряжения, скалярное и винтовое произведения, модуль и параметр бикватерниона. Представлены и другие объекты в пространстве на которые можно воздействовать винтовым движением: аффинная точки, вектор, прямая и плоскость, которые впоследствии можно представить в виде бикватернионов.

На основе принципа перенесения Котельникова–Штуди построены бикватернионные формулы винтового движения. Получены явные выражения для бикватерниона винтового преобразования, разложенного на вращательную и трансляционную составляющие. Выведены сэндвич-формулы для винтового движения точки, вектора, прямой и плоскости. Показана связь бикватернионного представления преобразований с матрицами проективных преобразований в трёхмерном пространстве. Дополнительно рассмотрены формулы отражения точки, плоскости и прямой относительно произвольной плоскости с помощью бикватернионов.

Осовные операции над бикватернионами и функции были реализованы в качестве блока на Julia. Блок содержит полный набор операций над бикватернионами и функций для описания пространственного перемещения объекта вокруг произвольной и фиксированной осей, а также отражения относительно произвольной плоскости. Использование Julia обеспечивает высокую производительность вычислений, что позволяет применять библиотеку в задачах, требующих численного моделирования в реальном времени.
Работоспособность программ проверена на тестовых примерах винтового движения точек. В качестве результата вычислений каждая программа выводит наглядный график.

Таким образом, разработанные программы удобны для вычисления и моделирования пространственных преобразований. 

